% ============================================================
%  GUÍAS DE IMPLEMENTACIÓN — COMUNICACIÓN EN SD
%  Sockets TCP/UDP · WebSockets · RPC · gRPC · RMI
%  Asignatura: Aplicaciones Distribuidas (ISR-701)
%  UTEQ — Ing. de Software — Nivel 7
%  Docente: Guerrero Ulloa Gleiston Ciceron
% ============================================================
\documentclass[12pt, a4paper, openany]{book}

% ── Encoding & fuentes ───────────────────────────────────────
\usepackage[utf8]{inputenc}
\usepackage[T1]{fontenc}
\usepackage[spanish,es-nodecimaldot]{babel}
\usepackage{csquotes}

% ── Márgenes y layout ────────────────────────────────────────
\usepackage[top=2.5cm, bottom=2.5cm, left=3cm, right=2.5cm]{geometry}
\usepackage{setspace}
\onehalfspacing

% ── Colores corporativos ─────────────────────────────────────
\usepackage[table,xcdraw,dvipsnames]{xcolor}
\definecolor{uteqGreen}    {RGB}{0,130,60}
\definecolor{uteqDark}     {RGB}{0,80,35}
\definecolor{uteqLight}    {RGB}{230,247,236}
\definecolor{practiceBlue} {RGB}{0,70,140}
\definecolor{practiceLight}{RGB}{230,240,255}
\definecolor{noteOrange}   {RGB}{180,83,9}
\definecolor{noteLight}    {RGB}{255,243,205}
\definecolor{warnRed}      {RGB}{185,28,28}
\definecolor{warnLight}    {RGB}{254,226,226}
\definecolor{codeGray}     {RGB}{246,248,250}
\definecolor{codeBorder}   {RGB}{180,190,200}
\definecolor{stepBlue}     {RGB}{37,99,235}

% ── Tipografía ───────────────────────────────────────────────
\usepackage{lmodern}
\usepackage{microtype}

% ── Matemáticas ──────────────────────────────────────────────
\usepackage{amsmath,amssymb,amsthm}

% ── Tablas ───────────────────────────────────────────────────
\usepackage{booktabs}
\usepackage{tabularx}
\usepackage{multirow}
\usepackage{array}

% ── Figuras ──────────────────────────────────────────────────
\usepackage{graphicx}
\usepackage{float}
\usepackage{caption}
\usepackage{subcaption}
\captionsetup{font=small, labelfont=bf}

% ── TikZ ─────────────────────────────────────────────────────
\usepackage{tikz}
\usetikzlibrary{arrows.meta, shapes, positioning, calc, fit,
	backgrounds, decorations.pathreplacing,
	decorations.markings, matrix, chains}
\tikzset{
	srvbox/.style={rectangle, rounded corners=4pt,
		draw=practiceBlue, fill=practiceLight,
		minimum width=2.0cm, minimum height=0.8cm,
		font=\small\bfseries, align=center},
	clibox/.style={rectangle, rounded corners=4pt,
		draw=uteqDark, fill=uteqLight,
		minimum width=2.0cm, minimum height=0.8cm,
		font=\small\bfseries, align=center},
	netbox/.style={rectangle, rounded corners=6pt,
		draw=gray!60, fill=gray!10,
		minimum width=1.8cm, minimum height=0.8cm,
		font=\small, align=center},
	msgArrow/.style={-Stealth, thick, draw=stepBlue},
	retArrow/.style={-Stealth, thick, draw=uteqDark, dashed},
	stepnum/.style={circle, draw=practiceBlue, fill=practiceBlue,
		text=white, font=\bfseries\small, minimum size=0.55cm, inner sep=0pt},
}

% ── Listings (código) ────────────────────────────────────────
\usepackage{listings}
\usepackage{lstautogobble}

% Caracteres especiales en listings
\lstset{
	inputencoding=utf8,
	extendedchars=true,
	literate=
	{á}{{\'{a}}}1 {é}{{\'{e}}}1 {í}{{\'{i}}}1
	{ó}{{\'{o}}}1 {ú}{{\'{u}}}1 {ñ}{{\~{n}}}1
	{Á}{{\'{A}}}1 {É}{{\'{E}}}1 {Í}{{\'{I}}}1
	{Ó}{{\'{O}}}1 {Ú}{{\'{U}}}1 {Ñ}{{\~{N}}}1
	{¿}{{\textquestiondown}}1 {¡}{{\textexclamdown}}1,
}

\lstdefinestyle{JavaStyle}{
	language=Java,
	basicstyle=\ttfamily\footnotesize,
	keywordstyle=\color{blue!80!black}\bfseries,
	commentstyle=\color{green!50!black}\itshape,
	stringstyle=\color{red!70!black},
	numberstyle=\tiny\color{gray},
	backgroundcolor=\color{codeGray},
	rulecolor=\color{codeBorder},
	numbers=left, numbersep=8pt, stepnumber=1,
	breaklines=true, breakatwhitespace=false,
	frame=single, framesep=5pt,
	tabsize=4, showstringspaces=false,
	captionpos=b, belowskip=0.8em, aboveskip=0.8em,
	mathescape=false, texcl=false,
	morekeywords={var, record, sealed, permits, yield, instanceof},
}

\lstdefinestyle{XMLStyle}{
	language=XML,
	basicstyle=\ttfamily\footnotesize,
	keywordstyle=\color{blue!70!black}\bfseries,
	stringstyle=\color{red!70!black},
	commentstyle=\color{green!50!black}\itshape,
	backgroundcolor=\color{noteLight},
	rulecolor=\color{noteOrange},
	numbers=left, numbersep=8pt,
	breaklines=true, frame=single, tabsize=2,
	captionpos=b, belowskip=0.8em, aboveskip=0.8em,
}

\lstdefinestyle{ProtoStyle}{
	basicstyle=\ttfamily\footnotesize,
	keywordstyle=\color{practiceBlue}\bfseries,
	stringstyle=\color{red!70!black},
	commentstyle=\color{green!50!black}\itshape,
	backgroundcolor=\color{practiceLight},
	rulecolor=\color{practiceBlue},
	numbers=left, numbersep=8pt,
	breaklines=true, frame=single, tabsize=2,
	captionpos=b, belowskip=0.8em, aboveskip=0.8em,
	morekeywords={syntax, package, service, message, rpc, returns,
		string, int32, int64, bool, float, double, repeated,
		stream, option, import},
}

\lstdefinestyle{BashStyle}{
	language=bash,
	basicstyle=\ttfamily\footnotesize\color{white},
	keywordstyle=\color{yellow!80!white}\bfseries,
	commentstyle=\color{gray!60!white}\itshape,
	stringstyle=\color{green!60!white},
	backgroundcolor=\color{black!85},
	rulecolor=\color{black},
	numbers=none, frame=single,
	breaklines=true, tabsize=2,
	captionpos=b, belowskip=0.8em, aboveskip=0.8em,
}

% ── tcolorbox ────────────────────────────────────────────────
\usepackage{tcolorbox}
\tcbuselibrary{skins, breakable, most}

\newtcolorbox{objetivobox}[1]{
	enhanced, breakable,
	colback=uteqLight, colframe=uteqGreen,
	fonttitle=\bfseries\color{white},
	title={Objetivo: #1},
	arc=4pt, boxrule=1pt,
	left=8pt, right=8pt, top=6pt, bottom=6pt
}

\newtcolorbox{notabox}{
	enhanced, breakable,
	colback=noteLight, colframe=noteOrange,
	fonttitle=\bfseries\color{white},
	title={Nota importante},
	arc=4pt, boxrule=1pt,
	left=8pt, right=8pt, top=6pt, bottom=6pt
}

\newtcolorbox{advertenciabox}{
	enhanced, breakable,
	colback=warnLight, colframe=warnRed,
	fonttitle=\bfseries\color{white},
	title={Advertencia},
	arc=4pt, boxrule=1pt,
	left=8pt, right=8pt, top=6pt, bottom=6pt
}

\newtcolorbox{prerequisitobox}{
	enhanced, breakable,
	colback=practiceLight, colframe=practiceBlue,
	fonttitle=\bfseries\color{white},
	title={Prerrequisitos},
	arc=4pt, boxrule=1pt,
	left=8pt, right=8pt, top=6pt, bottom=6pt
}

\newtcolorbox{resultadobox}{
	enhanced, breakable,
	colback=uteqLight, colframe=uteqGreen,
	fonttitle=\bfseries\color{white},
	title={Resultado esperado en consola},
	arc=4pt, boxrule=1pt,
	left=8pt, right=8pt, top=6pt, bottom=6pt
}

% ── Encabezados numerados de pasos ───────────────────────────
\newcounter{paso}[section]
\newcommand{\resetpasos}{\setcounter{paso}{0}}
\newcommand{\paso}[1]{%
	\stepcounter{paso}%
	\medskip\noindent
	\begin{tcolorbox}[enhanced, colback=white, colframe=practiceBlue,
		boxrule=0.8pt, arc=3pt, left=6pt, right=6pt, top=4pt, bottom=4pt,
		before upper={\noindent}]
		\textcolor{practiceBlue}{\textbf{Paso \arabic{paso}.}} #1
	\end{tcolorbox}%
}

% ── Hyperlinks ───────────────────────────────────────────────
\usepackage[colorlinks=true, linkcolor=practiceBlue,
urlcolor=uteqDark, citecolor=uteqGreen]{hyperref}
\usepackage{url}

% ── Bibliografía ─────────────────────────────────────────────
\usepackage[backend=biber, style=ieee, sorting=none]{biblatex}
\addbibresource{referencias_comunicacion.bib}

% Fix: listings Java language hace activos los chars < y > globalmente.
% Esto rompe TikZ. Restablecer catcodes al inicio del documento.
\AtBeginDocument{%
	\catcode`\>=12\relax
	\catcode`\<=12\relax
}

% ── Headers/Footers ──────────────────────────────────────────
\usepackage{fancyhdr}
\setlength{\headheight}{14pt}
\pagestyle{fancy}
\fancyhf{}
\fancyhead[L]{\small\textcolor{uteqDark}{\textbf{ISR-701 — Aplicaciones Distribuidas}}}
\fancyhead[R]{\small\textcolor{gray}{\leftmark}}
\fancyfoot[C]{\thepage}
\renewcommand{\headrulewidth}{0.4pt}

% ── Comandos auxiliares ───────────────────────────────────────
\newcommand{\code}[1]{\texttt{\small#1}}
\newcommand{\archivo}[1]{\texttt{\textbf{#1}}}
\newcommand{\paquete}[1]{\texttt{#1}}
\newcommand{\menu}[1]{\textit{#1}}

% ════════════════════════════════════════════════════════════
\begin{document}
	% ════════════════════════════════════════════════════════════
	
	% ── PORTADA ──────────────────────────────────────────────────
	\begin{titlepage}
		\centering
		{\large \textbf{Universidad Técnica Estatal de Quevedo}}\\[0.3em]
		{\normalsize Facultad de Ciencias de la Computación y Diseño Digital}\\
		{\normalsize Carrera de Ingeniería de Software}\\[2em]
		
		\rule{\linewidth}{1.5pt}\\[0.6em]
		{\LARGE \textbf{Guías de Implementación}}\\[0.4em]
		{\Large Mecanismos de Comunicación\\en Sistemas Distribuidos con Java}\\[0.4em]
		\rule{\linewidth}{1.5pt}\\[1.5em]
		
		{\large
			\begin{tabular}{rl}
				\textbf{Asignatura:}  & Aplicaciones Distribuidas (ISR-701)\\[0.3em]
				\textbf{Unidad:}      & 1 — Generalidades de los Sistemas Distribuidos\\[0.3em]
				\textbf{Nivel:}       & Séptimo semestre — Ingeniería de Software\\[0.3em]
				\textbf{IDE:}         & IntelliJ IDEA 2024.x\\[0.3em]
				\textbf{Lenguaje:}    & Java 21 (LTS) + Maven 3.9+\\[0.3em]
				\textbf{Período:}     & Regular 2026-2027\\[0.3em]
				\textbf{Docente:}     & Ing. Gleiston Guerrero Ulloa, Mg.Sc.\\
			\end{tabular}
		}\\[2em]
		
		\begin{tcolorbox}[colback=uteqLight, colframe=uteqGreen,
			fonttitle=\bfseries\color{white}, title={Contenido de este documento},
			arc=4pt, boxrule=1pt]
			\begin{enumerate}
				\item \textbf{Guía 1} — Sockets TCP y UDP
				\item \textbf{Guía 2} — WebSockets con Jakarta EE / Spring Boot
				\item \textbf{Guía 3} — RPC con Apache Thrift
				\item \textbf{Guía 4} — gRPC con Protocol Buffers
				\item \textbf{Guía 5} — RMI (Remote Method Invocation)
			\end{enumerate}
			Cada guía incluye: fundamento teórico, diagrama de arquitectura,
			código fuente completo paso a paso, ejecución desde IntelliJ IDEA
			y verificación de resultados.
		\end{tcolorbox}
		
		\vfill
		{\small UTEQ · Quevedo, Los Ríos, Ecuador · \today}
	\end{titlepage}
	
	% ── TOC ──────────────────────────────────────────────────────
	\tableofcontents
	\listoffigures
	\newpage
	
	% ════════════════════════════════════════════════════════════
	%  INTRODUCCIÓN TRANSVERSAL
	% ════════════════════════════════════════════════════════════
	\chapter*{Introducción}
	\addcontentsline{toc}{chapter}{Introducción}
	\markboth{Introducción}{Introducción}
	
	Los sistemas distribuidos modernos utilizan múltiples mecanismos de comunicación
	entre procesos (\emph{Inter-Process Communication}, IPC). La elección del mecanismo
	correcto determina la latencia, el throughput, la complejidad del código y la
	interoperabilidad del sistema~\parencite{vanSteen2023}.
	
	Este documento presenta cinco guías de implementación progresivas para Java~21 con
	IntelliJ IDEA, cubriendo el espectro completo desde los sockets de bajo nivel hasta
	la invocación de métodos remotos de alto nivel:
	
	\begin{table}[H]
		\centering
		\caption{Comparativa de mecanismos de comunicación cubiertos en las guías}
		\label{tab:comparativa}
		\begin{tabularx}{\linewidth}{@{}lXXXX@{}}
			\toprule
			\textbf{Mecanismo} & \textbf{Nivel} & \textbf{Protocolo} & \textbf{Tipado} & \textbf{Uso típico} \\
			\midrule
			Sockets TCP  & Bajo   & TCP/IP   & Sin tipado & Streaming, chat \\
			Sockets UDP  & Bajo   & UDP/IP   & Sin tipado & VoIP, gaming, DNS \\
			WebSocket    & Medio  & HTTP/WS  & Sin tipado & Tiempo real, notif. \\
			RPC (Thrift) & Alto   & TCP/bin  & Fuerte     & Microservicios \\
			gRPC         & Alto   & HTTP/2   & Fuerte     & Microservicios \\
			RMI          & Alto   & JRMP     & Java       & Sistemas Java \\
			\bottomrule
		\end{tabularx}
	\end{table}
	
	\begin{prerequisitobox}
		Antes de comenzar cualquier guía, verificar la instalación de:
		\begin{itemize}
			\item \textbf{JDK 21} (Adoptium Temurin 21 LTS): \code{java -\/\/-version}
			\item \textbf{Maven 3.9+}: \code{mvn -\/\/-version}
			\item \textbf{IntelliJ IDEA} 2024.x (Community o Ultimate)
			\item \textbf{Git} 2.x: para control de versiones del código
		\end{itemize}
	\end{prerequisitobox}
	
	% ════════════════════════════════════════════════════════════
	%  GUÍA 1 — SOCKETS TCP/UDP
	% ════════════════════════════════════════════════════════════
	\chapter{Guía 1: Sockets TCP y UDP}
	\resetpasos
	
	\begin{objetivobox}{Sockets TCP y UDP}
		Implementar un sistema de mensajería distribuida usando sockets TCP
		(orientado a conexión, entrega garantizada) y UDP (sin conexión, baja latencia),
		comprendiendo las diferencias arquitectónicas entre ambos protocolos~\parencite{stevens2004}.
	\end{objetivobox}
	
	\section{Fundamento teórico}
	
	Un \textbf{socket} es la abstracción del sistema operativo que permite a dos procesos
	en red intercambiar datos. Representa un endpoint de comunicación identificado por la
	tupla \texttt{(IP, puerto, protocolo)}~\parencite{coulouris2012}.
	
	\subsection{TCP vs. UDP}
	
	\textbf{TCP} (\textit{Transmission Control Protocol}) establece una conexión antes de
	transferir datos mediante el \emph{three-way handshake} (SYN $\to$ SYN-ACK $\to$ ACK).
	Garantiza: entrega en orden, detección y retransmisión de pérdidas, y control de flujo
	mediante ventana deslizante~\parencite{kurose2017}.
	
	\textbf{UDP} (\textit{User Datagram Protocol}) envía datagramas independientes sin
	establecer conexión previa. Ofrece latencia mínima al costo de no garantizar orden ni
	entrega, siendo adecuado cuando la pérdida ocasional de paquetes es tolerable
	(VoIP, juegos en línea, DNS)~\parencite{kurose2017}.
	
	\begin{figure}[H]
		\centering
		\begin{tikzpicture}[node distance=1.0cm and 2.5cm, font=\small]
			% TCP
			\node[clibox] (cliTCP) {Cliente TCP};
			\node[srvbox, right=3.0cm of cliTCP] (srvTCP) {Servidor TCP};
			\draw[msgArrow] (cliTCP.north east) -- node[above, font=\tiny]{SYN}
			([yshift=4pt]srvTCP.north west);
			\draw[retArrow] ([yshift=-4pt]srvTCP.north west) -- node[below, font=\tiny]{SYN-ACK}
			(cliTCP.south east);
			\draw[msgArrow] ([yshift=8pt]cliTCP.east) -- node[above, font=\tiny]{ACK}
			([yshift=8pt]srvTCP.west);
			\draw[msgArrow] (cliTCP.east) -- node[above, font=\tiny]{datos}
			(srvTCP.west);
			% UDP
			\node[clibox, below=2.0cm of cliTCP] (cliUDP) {Cliente UDP};
			\node[srvbox, right=3.0cm of cliUDP] (srvUDP) {Servidor UDP};
			\draw[msgArrow] (cliUDP.east) -- node[above, font=\tiny]{datagrama (sin handshake)}
			(srvUDP.west);
			\node[left=0.2cm of cliTCP, font=\bfseries\small, text=practiceBlue] {TCP};
			\node[left=0.2cm of cliUDP, font=\bfseries\small, text=uteqDark] {UDP};
		\end{tikzpicture}
		\caption{Comparativa del flujo de establecimiento de conexión TCP vs. UDP}
		\label{fig:tcp_udp}
	\end{figure}
	
	\section{Estructura del proyecto Maven}
	
	\paso{Abrir IntelliJ IDEA. Ir a \menu{File $\to$ New $\to$ Project}.
		Seleccionar \textbf{Maven Archetype}, elegir \code{org.apache.maven.archetypes:maven-archetype-quickstart},
		definir:}
	
	\begin{itemize}
		\item \textbf{GroupId:} \code{ec.edu.uteq.distribuidas}
		\item \textbf{ArtifactId:} \code{guia01-sockets}
		\item \textbf{Version:} \code{1.0-SNAPSHOT}
	\end{itemize}
	
	\paso{Reemplazar el contenido de \archivo{pom.xml} con la siguiente configuración
		para Java 21:}
	
	\begin{lstlisting}[style=XMLStyle, caption={pom.xml --- Configuración del proyecto}]
		<?xml version="1.0" encoding="UTF-8"?>
		<project xmlns="http://maven.apache.org/POM/4.0.0"
		xmlns:xsi="http://www.w3.org/2001/XMLSchema-instance"
		xsi:schemaLocation="http://maven.apache.org/POM/4.0.0
		http://maven.apache.org/xsd/maven-4.0.0.xsd">
		<modelVersion>4.0.0</modelVersion>
		<groupId>ec.edu.uteq.distribuidas</groupId>
		<artifactId>guia01-sockets</artifactId>
		<version>1.0-SNAPSHOT</version>
		<properties>
		<maven.compiler.source>21</maven.compiler.source>
		<maven.compiler.target>21</maven.compiler.target>
		<project.build.sourceEncoding>UTF-8</project.build.sourceEncoding>
		</properties>
		</project>
	\end{lstlisting}
	
	\section{Implementación de Sockets TCP}
	
	\paso{Crear el paquete \code{ec.edu.uteq.distribuidas.tcp} dentro de
		\code{src/main/java}. Clic derecho sobre \code{ec.edu.uteq.distribuidas}
		$\to$ \menu{New $\to$ Package} $\to$ ingresar \code{tcp}.}
	
	\paso{Crear la clase \archivo{ServidorTCP.java} en el paquete \code{tcp}:}
	
	\begin{lstlisting}[style=JavaStyle, caption={ServidorTCP.java --- Servidor multihilo con ThreadPool}]
		package ec.edu.uteq.distribuidas.tcp;
		
		import java.io.*;
		import java.net.*;
		import java.time.LocalDateTime;
		import java.time.format.DateTimeFormatter;
		import java.util.concurrent.*;
		import java.util.concurrent.atomic.*;
		
		/**
		* Servidor TCP multihilo con pool de hilos fijo.
		* Escucha en el puerto 9000 y atiende comandos:
		*   HORA   -> retorna la hora actual del servidor
		*   ECO    -> retorna el mensaje prefijado con ECO:
		*   SALIR  -> cierra la conexion del cliente
		*/
		public class ServidorTCP {
			
			private static final int PUERTO = 9000;
			private static final int MAX_HILOS = 10;
			private static final AtomicInteger contadorClientes = new AtomicInteger(0);
			
			public static void main(String[] args) throws IOException {
				ExecutorService pool = Executors.newFixedThreadPool(MAX_HILOS);
				
				try (ServerSocket servidor = new ServerSocket(PUERTO)) {
					System.out.println("[" + timestamp() + "] Servidor listo en :" + PUERTO);
					
					while (true) {
						Socket cliente = servidor.accept();   // BLOQUEA hasta nueva conexion
						int id = contadorClientes.incrementAndGet();
						System.out.println("[" + timestamp() + "] Cliente #" + id
						+ " conectado desde " + cliente.getInetAddress());
						pool.submit(new ManejadorCliente(cliente, id));
					}
				}
			}
			
			public static String timestamp() {
				return LocalDateTime.now()
				.format(DateTimeFormatter.ofPattern("HH:mm:ss.SSS"));
			}
			
			// Clase interna para manejar cada cliente en su propio hilo
			static class ManejadorCliente implements Runnable {
				private final Socket socket;
				private final int idCliente;
				
				ManejadorCliente(Socket socket, int idCliente) {
					this.socket   = socket;
					this.idCliente = idCliente;
				}
				
				@Override
				public void run() {
					try (BufferedReader entrada  = new BufferedReader(
					new InputStreamReader(socket.getInputStream()));
					PrintWriter   salida   = new PrintWriter(
					socket.getOutputStream(), true)) {
						
						String linea;
						while ((linea = entrada.readLine()) != null) {
							System.out.println("[" + ServidorTCP.timestamp()
							+ "] Cliente#" + idCliente + " > " + linea);
							
							if (linea.equalsIgnoreCase("HORA")) {
								salida.println("HORA_ACTUAL:" + ServidorTCP.timestamp());
								
							} else if (linea.toUpperCase().startsWith("ECO")) {
								String mensaje = linea.length() > 4
								? linea.substring(4) : "(vacio)";
								salida.println("ECO[#" + idCliente + "]:" + mensaje);
								
							} else if (linea.equalsIgnoreCase("SALIR")) {
								salida.println("ADIOS:hasta luego cliente #" + idCliente);
								System.out.println("[" + ServidorTCP.timestamp()
								+ "] Cliente #" + idCliente + " desconectado.");
								break;
								
							} else {
								salida.println("ERROR:comando desconocido -> " + linea);
							}
						}
					} catch (IOException e) {
						System.err.println("Error cliente #" + idCliente + ": " + e.getMessage());
					}
				}
			}
		}
	\end{lstlisting}
	
	\paso{Crear la clase \archivo{ClienteTCP.java} en el mismo paquete \code{tcp}:}
	
	\begin{lstlisting}[style=JavaStyle, caption={ClienteTCP.java --- Cliente interactivo}]
		package ec.edu.uteq.distribuidas.tcp;
		
		import java.io.*;
		import java.net.*;
		import java.util.Scanner;
		
		/**
		* Cliente TCP interactivo.
		* Envía comandos al servidor y muestra la respuesta.
		* Uso: ejecutar desde IntelliJ -> Run 'ClienteTCP.main()'
		*/
		public class ClienteTCP {
			
			private static final String HOST  = "localhost";
			private static final int    PUERTO = 9000;
			
			public static void main(String[] args) throws IOException {
				try (Socket socket  = new Socket(HOST, PUERTO);
				BufferedReader entrada = new BufferedReader(
				new InputStreamReader(socket.getInputStream()));
				PrintWriter salida = new PrintWriter(
				socket.getOutputStream(), true);
				Scanner teclado = new Scanner(System.in)) {
					
					System.out.println("Conectado a " + HOST + ":" + PUERTO);
					System.out.println("Comandos disponibles: HORA | ECO <texto> | SALIR");
					System.out.println("---------------------------------------------------");
					
					String comando;
					while (true) {
						System.out.print("> ");
						comando = teclado.nextLine().trim();
						if (comando.isEmpty()) continue;
						
						salida.println(comando);   // enviar al servidor
						String respuesta = entrada.readLine();
						System.out.println("Servidor: " + respuesta);
						
						if (comando.equalsIgnoreCase("SALIR")) break;
					}
				}
				System.out.println("Conexion cerrada.");
			}
		}
	\end{lstlisting}
	
	\section{Implementación de Sockets UDP}
	
	\paso{Crear el paquete \code{ec.edu.uteq.distribuidas.udp} y la clase
		\archivo{ServidorUDP.java}:}
	
	\begin{lstlisting}[style=JavaStyle, caption={ServidorUDP.java --- Servidor sin conexión}]
		package ec.edu.uteq.distribuidas.udp;
		
		import java.net.*;
		import java.time.LocalDateTime;
		import java.time.format.DateTimeFormatter;
		
		/**
		* Servidor UDP sin conexion.
		* Recibe datagramas de hasta 1024 bytes y responde en el
		* mismo socket. No tiene estado entre paquetes.
		*/
		public class ServidorUDP {
			
			private static final int PUERTO   = 9001;
			private static final int TAM_BUF  = 1024;
			
			public static void main(String[] args) throws Exception {
				try (DatagramSocket socket = new DatagramSocket(PUERTO)) {
					System.out.println("Servidor UDP escuchando en puerto " + PUERTO);
					
					byte[] buffer = new byte[TAM_BUF];
					while (true) {
						DatagramPacket paquete =
						new DatagramPacket(buffer, buffer.length);
						socket.receive(paquete);   // BLOQUEA hasta recibir datagrama
						
						String mensaje = new String(paquete.getData(),
						0, paquete.getLength());
						String hora = LocalDateTime.now()
						.format(DateTimeFormatter.ofPattern("HH:mm:ss.SSS"));
						
						System.out.println("[" + hora + "] Recibido de "
						+ paquete.getAddress() + ":" + paquete.getPort()
						+ " -> " + mensaje);
						
						// Preparar respuesta
						String respuesta = "ACK:" + mensaje + " @" + hora;
						byte[] resp = respuesta.getBytes();
						DatagramPacket respPkt = new DatagramPacket(
						resp, resp.length,
						paquete.getAddress(), paquete.getPort());
						socket.send(respPkt);
					}
				}
			}
		}
	\end{lstlisting}
	
	\begin{lstlisting}[style=JavaStyle, caption={ClienteUDP.java --- Cliente sin conexión}]
		package ec.edu.uteq.distribuidas.udp;
		
		import java.net.*;
		import java.util.Scanner;
		
		/**
		* Cliente UDP: envía un datagrama y espera respuesta con
		* timeout de 3 segundos.
		*/
		public class ClienteUDP {
			
			private static final String HOST    = "localhost";
			private static final int    PUERTO  = 9001;
			private static final int    TIMEOUT = 3000;  // ms
			private static final int    TAM_BUF = 1024;
			
			public static void main(String[] args) throws Exception {
				try (DatagramSocket socket = new DatagramSocket()) {
					socket.setSoTimeout(TIMEOUT);
					InetAddress servidor = InetAddress.getByName(HOST);
					Scanner sc = new Scanner(System.in);
					
					System.out.println("Cliente UDP. Escribe un mensaje (o 'salir'):");
					while (sc.hasNextLine()) {
						String linea = sc.nextLine().trim();
						if (linea.equalsIgnoreCase("salir")) break;
						
						byte[] datos = linea.getBytes();
						DatagramPacket pkt = new DatagramPacket(
						datos, datos.length, servidor, PUERTO);
						socket.send(pkt);
						
						// Esperar respuesta
						byte[] buf = new byte[TAM_BUF];
						DatagramPacket resp = new DatagramPacket(buf, buf.length);
						try {
							socket.receive(resp);
							System.out.println("Respuesta: "
							+ new String(resp.getData(), 0, resp.getLength()));
						} catch (SocketTimeoutException e) {
							System.out.println("Timeout: no se recibio respuesta en "
							+ TIMEOUT + "ms");
						}
					}
				}
			}
		}
	\end{lstlisting}
	
	\section{Ejecución desde IntelliJ IDEA}
	
	\paso{Ejecutar \textbf{ServidorTCP}: clic derecho sobre la clase $\to$
		\menu{Run 'ServidorTCP.main()'}. Debe aparecer:}
	
	\begin{resultadobox}
		\begin{lstlisting}[style=BashStyle]
			[10:05:32.001] Servidor listo en :9000
		\end{lstlisting}
	\end{resultadobox}
	
	\paso{Ejecutar \textbf{ClienteTCP}: ir a \menu{Run $\to$ Edit Configurations
			$\to$ + $\to$ Application}, configurar la clase \code{ClienteTCP} y ejecutar.
		Probar los comandos:}
	
	\begin{resultadobox}
		\begin{lstlisting}[style=BashStyle]
			> HORA
			Servidor: HORA_ACTUAL:10:05:40.234
			> ECO Hola UTEQ
			Servidor: ECO[#1]:Hola UTEQ
			> SALIR
			Servidor: ADIOS:hasta luego cliente #1
		\end{lstlisting}
	\end{resultadobox}
	
	\paso{Para el UDP: ejecutar \textbf{ServidorUDP} primero, luego
		\textbf{ClienteUDP}. Para ejecutar múltiples instancias simultáneas del cliente,
		activar \menu{Allow multiple instances} en la configuración de ejecución.}
	
	\begin{notabox}
		TCP usa \code{ServerSocket.accept()} que bloquea hasta recibir una conexión.
		Cada cliente se atiende en un hilo separado del \code{ExecutorService}.
		UDP usa \code{DatagramSocket.receive()} que bloquea hasta recibir un datagrama,
		sin estado entre paquetes.
	\end{notabox}
	
	% ════════════════════════════════════════════════════════════
	%  GUÍA 2 — WEBSOCKETS
	% ════════════════════════════════════════════════════════════
	\chapter{Guía 2: WebSockets con Spring Boot}
	\resetpasos
	
	\begin{objetivobox}{WebSockets con Spring Boot}
		Implementar un servidor de chat en tiempo real usando el protocolo WebSocket
		sobre HTTP/1.1, habilitando comunicación bidireccional y persistente entre
		clientes y servidor~\parencite{fette2011}.
	\end{objetivobox}
	
	\section{Fundamento teórico}
	
	\textbf{WebSocket} (RFC 6455~\parencite{fette2011}) es un protocolo de comunicación
	que establece un canal full-duplex sobre una única conexión TCP. A diferencia de
	HTTP, que es petición-respuesta (half-duplex), WebSocket mantiene la conexión
	abierta permitiendo que el servidor envíe datos al cliente sin que este los solicite.
	
	El ciclo de vida tiene tres fases:
	\begin{enumerate}
		\item \textbf{Handshake HTTP}: el cliente envía una cabecera
		\code{Upgrade: websocket}. El servidor responde \code{101 Switching Protocols}.
		\item \textbf{Transferencia de frames}: datos se transmiten en frames binarios
		o de texto, eficientemente.
		\item \textbf{Cierre}: cualquiera de los extremos envía un frame \code{CLOSE}.
	\end{enumerate}
	
	\begin{figure}[H]
		\centering
		\begin{tikzpicture}[node distance=0.5cm and 3.5cm, font=\small]
			\node[clibox] (cli) {Browser / Cliente};
			\node[srvbox, right=4.0cm of cli] (srv) {Servidor Spring Boot};
			
			\draw[msgArrow] ([yshift=30pt]cli.east)  --
			node[above, font=\tiny]{GET /ws HTTP/1.1\quad Upgrade: websocket}
			([yshift=30pt]srv.west);
			\draw[retArrow] ([yshift=18pt]srv.west) --
			node[below, font=\tiny]{101 Switching Protocols}
			([yshift=18pt]cli.east);
			\draw[msgArrow] ([yshift=4pt]cli.east)  --
			node[above, font=\tiny]{frame texto: ''Hola sala''}
			([yshift=4pt]srv.west);
			\draw[retArrow] ([yshift=-8pt]srv.west) --
			node[below, font=\tiny]{frame texto: ''[user] Hola sala'' (broadcast)}
			([yshift=-8pt]cli.east);
			\draw[msgArrow] ([yshift=-22pt]cli.east) --
			node[above, font=\tiny]{frame CLOSE}
			([yshift=-22pt]srv.west);
		\end{tikzpicture}
		\caption{Ciclo de vida de una conexión WebSocket sobre HTTP}
		\label{fig:websocket_lifecycle}
	\end{figure}
	
	\section{Crear proyecto Spring Boot}
	
	\paso{Abrir IntelliJ IDEA $\to$ \menu{File $\to$ New $\to$ Project $\to$
			Spring Initializr}. Configurar:}
	
	\begin{itemize}
		\item \textbf{Name:} \code{guia02-websocket}
		\item \textbf{Language:} Java
		\item \textbf{Type:} Maven
		\item \textbf{Java:} 21
		\item \textbf{Package name:} \code{ec.edu.uteq.distribuidas.ws}
	\end{itemize}
	
	\paso{En la pantalla de dependencias, buscar y agregar:
		\textbf{WebSocket} (Spring WebSocket). Hacer clic en \menu{Finish}.}
	
	El \archivo{pom.xml} generado incluirá automáticamente:
	
	\begin{lstlisting}[style=XMLStyle, caption={Dependencia WebSocket en pom.xml}]
		<dependency>
		<groupId>org.springframework.boot</groupId>
		<artifactId>spring-boot-starter-websocket</artifactId>
		</dependency>
		<dependency>
		<groupId>org.webjars</groupId>
		<artifactId>webjars-locator-core</artifactId>
		</dependency>
	\end{lstlisting}
	
	\section{Configuración del servidor WebSocket}
	
	\paso{Crear la clase de configuración \archivo{WebSocketConfig.java}:}
	
	\begin{lstlisting}[style=JavaStyle, caption={WebSocketConfig.java --- Configuración STOMP sobre WebSocket}]
		package ec.edu.uteq.distribuidas.ws.config;
		
		import org.springframework.context.annotation.Configuration;
		import org.springframework.messaging.simp.config.MessageBrokerRegistry;
		import org.springframework.web.socket.config.annotation.*;
		
		/**
		* Configura el broker de mensajes STOMP sobre WebSocket.
		* STOMP (Simple Text Oriented Messaging Protocol) agrega
		* un protocolo de mensajeria sobre el canal WebSocket raw.
		*/
		@Configuration
		@EnableWebSocketMessageBroker
		public class WebSocketConfig implements WebSocketMessageBrokerConfigurer {
			
			@Override
			public void configureMessageBroker(MessageBrokerRegistry config) {
				// Prefijo de los topics a los que los clientes se suscriben
				config.enableSimpleBroker("/topic");
				// Prefijo de los mensajes enviados al servidor
				config.setApplicationDestinationPrefixes("/app");
			}
			
			@Override
			public void registerStompEndpoints(StompEndpointRegistry registry) {
				// Endpoint WebSocket con fallback SockJS para navegadores sin soporte
				registry.addEndpoint("/ws-chat")
				.setAllowedOriginPatterns("*")
				.withSockJS();
			}
		}
	\end{lstlisting}
	
	\paso{Crear el modelo de mensaje \archivo{MensajeChat.java}:}
	
	\begin{lstlisting}[style=JavaStyle, caption={MensajeChat.java --- Modelo de mensaje del chat}]
		package ec.edu.uteq.distribuidas.ws.model;
		
		import java.time.LocalDateTime;
		import java.time.format.DateTimeFormatter;
		
		/**
		* Record inmutable que representa un mensaje de chat.
		* Los records de Java 16+ generan automaticamente constructor,
		* getters, equals, hashCode y toString.
		*/
		public record MensajeChat(
		String usuario,
		String contenido,
		String tipo,     // "UNIRSE" | "MENSAJE" | "SALIR"
		String timestamp
		) {
			// Constructor compacto: valida y asigna timestamp si es null
			public MensajeChat {
				if (timestamp == null) {
					timestamp = LocalDateTime.now()
					.format(DateTimeFormatter.ofPattern("HH:mm:ss"));
				}
			}
		}
	\end{lstlisting}
	
	\paso{Crear el controlador de mensajes \archivo{ChatController.java}:}
	
	\begin{lstlisting}[style=JavaStyle, caption={ChatController.java --- Manejador de mensajes STOMP}]
		package ec.edu.uteq.distribuidas.ws.controller;
		
		import ec.edu.uteq.distribuidas.ws.model.MensajeChat;
		import org.springframework.messaging.handler.annotation.*;
		import org.springframework.messaging.simp.SimpMessageSendingOperations;
		import org.springframework.stereotype.Controller;
		
		/**
		* Controlador de mensajes WebSocket.
		* @MessageMapping mapea los mensajes STOMP enviados al servidor.
		* @SendTo especifica el topic al que se reenvian los mensajes.
		*/
		@Controller
		public class ChatController {
			
			private final SimpMessageSendingOperations messagingTemplate;
			
			public ChatController(SimpMessageSendingOperations messagingTemplate) {
				this.messagingTemplate = messagingTemplate;
			}
			
			/**
			* Recibe mensajes en /app/chat.enviar
			* y los reenvía a todos los suscriptores de /topic/sala-general.
			*/
			@MessageMapping("/chat.enviar")
			@SendTo("/topic/sala-general")
			public MensajeChat enviarMensaje(MensajeChat mensaje) {
				System.out.println("[WS] " + mensaje.usuario() + ": " + mensaje.contenido());
				return mensaje;
			}
			
			/**
			* Notifica a toda la sala cuando un usuario nuevo se une.
			*/
			@MessageMapping("/chat.unirse")
			@SendTo("/topic/sala-general")
			public MensajeChat unirse(MensajeChat mensaje,
			@Header("simpSessionId") String sessionId) {
				System.out.println("[WS] Usuario '" + mensaje.usuario()
				+ "' entró. SessionId: " + sessionId);
				return new MensajeChat(mensaje.usuario(),
				mensaje.usuario() + " se unió a la sala.",
				"UNIRSE", null);
			}
		}
	\end{lstlisting}
	
	\paso{Crear el cliente HTML \archivo{chat.html} en \code{src/main/resources/static}
		para probar desde el navegador:}
	
	\begin{lstlisting}[style=XMLStyle, caption={chat.html --- Cliente WebSocket en el navegador}]
		<!DOCTYPE html>
		<html lang="es">
		<head>
		<meta charset="UTF-8"/>
		<title>Chat WebSocket - UTEQ</title>
		<script src="https://cdn.jsdelivr.net/npm/sockjs-client@1/dist/sockjs.min.js"></script>
		<script src="https://cdn.jsdelivr.net/npm/stompjs@2.3.3/lib/stomp.min.js"></script>
		</head>
		<body>
		<h2>Chat en Tiempo Real (WebSocket + STOMP)</h2>
		<input id="usuario" placeholder="Tu nombre" />
		<button onclick="conectar()">Conectar</button>
		<hr/>
		<div id="mensajes" style="height:300px;overflow-y:auto;border:1px solid #ccc"></div>
		<input id="texto" placeholder="Escribe un mensaje..." style="width:70%"/>
		<button onclick="enviar()">Enviar</button>
		
		<script>
		let stompClient = null;
		
		function conectar() {
			const socket = new SockJS('/ws-chat');
			stompClient = Stomp.over(socket);
			stompClient.connect({}, frame => {
				// Suscribirse al topic de la sala
				stompClient.subscribe('/topic/sala-general', msg => {
					const m = JSON.parse(msg.body);
					mostrar('[' + m.timestamp + '] ' + m.usuario + ': ' + m.contenido);
				});
				// Anunciar ingreso
				stompClient.send('/app/chat.unirse', {},
				JSON.stringify({usuario: document.getElementById('usuario').value,
					contenido: '', tipo: 'UNIRSE', timestamp: null}));
			});
		}
		
		function enviar() {
			const texto = document.getElementById('texto').value;
			stompClient.send('/app/chat.enviar', {},
			JSON.stringify({usuario: document.getElementById('usuario').value,
				contenido: texto, tipo: 'MENSAJE', timestamp: null}));
			document.getElementById('texto').value = '';
		}
		
		function mostrar(txt) {
			const div = document.getElementById('mensajes');
			div.innerHTML += '<p>' + txt + '</p>';
			div.scrollTop = div.scrollHeight;
		}
		</script>
		</body>
		</html>
	\end{lstlisting}
	
	\section{Ejecución}
	
	\paso{Ejecutar la clase principal de Spring Boot: clic derecho $\to$
		\menu{Run 'Guia02WebsocketApplication'}. Esperar hasta ver:}
	
	\begin{resultadobox}
		\begin{lstlisting}[style=BashStyle]
			Started Guia02WebsocketApplication in 2.3s (JVM running for 2.8)
			Tomcat started on port 8080 (http)
		\end{lstlisting}
	\end{resultadobox}
	
	\paso{Abrir \textbf{dos} pestañas de navegador en
		\url{http://localhost:8080/chat.html}. En cada pestaña introducir
		un nombre de usuario diferente, hacer clic en \textbf{Conectar} y enviar
		mensajes. Los mensajes deben aparecer en \emph{ambas} pestañas en tiempo real.}
	
	% ════════════════════════════════════════════════════════════
	%  GUÍA 3 — RPC CON APACHE THRIFT
	% ════════════════════════════════════════════════════════════
	\chapter{Guía 3: RPC con Apache Thrift}
	\resetpasos
	
	\begin{objetivobox}{RPC con Apache Thrift}
		Implementar un sistema de llamada a procedimiento remoto (RPC) usando Apache
		Thrift, que genera automáticamente el código de serialización (stub/skeleton)
		a partir de un archivo de definición de interfaz (\code{.thrift})~\parencite{slee2007}.
	\end{objetivobox}
	
	\section{Fundamento teórico}
	
	\textbf{RPC} (\textit{Remote Procedure Call}) permite a un proceso invocar un
	procedimiento en otro proceso, posiblemente en otra máquina, como si fuera local.
	El mecanismo transparente de comunicación se logra mediante dos componentes
	generados automáticamente~\parencite{birrell1984}:
	
	\begin{itemize}
		\item \textbf{Stub cliente} (proxy): intercepta la llamada, serializa los
		parámetros (\emph{marshalling}), los envía por la red y devuelve el resultado.
		\item \textbf{Skeleton servidor}: recibe los bytes, los deserializa
		(\emph{unmarshalling}) e invoca la implementación real.
	\end{itemize}
	
	\textbf{Apache Thrift}~\parencite{slee2007} es un framework de Facebook (2007) que
	define interfaces en un IDL (\textit{Interface Definition Language}) propio y genera
	código para más de 20 lenguajes.
	
	\section{Configuración del proyecto}
	
	\paso{Crear un proyecto Maven en IntelliJ con GroupId
		\code{ec.edu.uteq.distribuidas}, ArtifactId \code{guia03-rpc}.}
	
	\paso{Agregar la dependencia de Thrift en \archivo{pom.xml}:}
	
	\begin{lstlisting}[style=XMLStyle, caption={pom.xml --- Dependencia Apache Thrift}]
		<dependencies>
		<dependency>
		<groupId>org.apache.thrift</groupId>
		<artifactId>libthrift</artifactId>
		<version>0.19.0</version>
		</dependency>
		</dependencies>
		
		<build>
		<plugins>
		<plugin>
		<groupId>org.apache.thrift.tools</groupId>
		<artifactId>maven-thrift-plugin</artifactId>
		<version>0.1.11</version>
		<configuration>
		<thriftExecutable>thrift</thriftExecutable>
		<generator>java</generator>
		</configuration>
		<executions>
		<execution>
		<id>thrift-sources</id>
		<phase>generate-sources</phase>
		<goals><goal>compile</goal></goals>
		</execution>
		</executions>
		</plugin>
		</plugins>
		</build>
	\end{lstlisting}
	
	\begin{notabox}
		Si Thrift no está instalado en el sistema, descargarlo de
		\url{https://thrift.apache.org/download} y agregar el ejecutable al PATH.
		En Ubuntu/Debian: \code{sudo apt install thrift-compiler}.
	\end{notabox}
	
	\section{Definición del IDL}
	
	\paso{Crear el directorio \code{src/main/thrift} y dentro el archivo
		\archivo{calculadora.thrift}:}
	
	\begin{lstlisting}[style=ProtoStyle, caption={calculadora.thrift --- Definición IDL de la calculadora remota}]
		/**
		* Definicion IDL de Thrift para una Calculadora Remota.
		* El compilador de Thrift genera automaticamente:
		*   - Calculadora.java (interfaz del servicio)
		*   - Calculadora.Client (stub del cliente)
		*   - Calculadora.Processor (skeleton del servidor)
		*/
		namespace java ec.edu.uteq.distribuidas.rpc.gen
		
		/** Excepcion de dominio para operaciones invalidas */
		exception OperacionInvalidaException {
			1: required string mensaje
		}
		
		/** Servicio de calculadora distribuida */
		service Calculadora {
			
			/** Suma dos numeros de doble precision */
			double sumar(1: double a, 2: double b)
			
			/** Resta b de a */
			double restar(1: double a, 2: double b)
			
			/** Multiplica dos numeros */
			double multiplicar(1: double a, 2: double b)
			
			/** Divide a entre b. Lanza excepcion si b == 0 */
			double dividir(1: double a, 2: double b)
			throws (1: OperacionInvalidaException ex)
			
			/** Retorna el nombre del servidor que atendio */
			string nombreServidor()
		}
	\end{lstlisting}
	
	\paso{Generar el código Java. En la terminal integrada de IntelliJ
		(\menu{View $\to$ Tool Windows $\to$ Terminal}):}
	
	\begin{lstlisting}[style=BashStyle]
		# Desde la raiz del proyecto
		thrift --gen java -out src/main/java src/main/thrift/calculadora.thrift
	\end{lstlisting}
	
	Esto genera en \code{src/main/java/ec/edu/uteq/distribuidas/rpc/gen/}
	los archivos: \archivo{Calculadora.java} y \archivo{OperacionInvalidaException.java}.
	
	\section{Implementación servidor y cliente}
	
	\paso{Implementar la lógica del servidor en \archivo{CalculadoraImpl.java}:}
	
	\begin{lstlisting}[style=JavaStyle, caption={CalculadoraImpl.java --- Implementación de la lógica RPC}]
		package ec.edu.uteq.distribuidas.rpc.server;
		
		import ec.edu.uteq.distribuidas.rpc.gen.Calculadora;
		import ec.edu.uteq.distribuidas.rpc.gen.OperacionInvalidaException;
		import org.apache.thrift.TException;
		
		/**
		* Implementación concreta de la interfaz generada por Thrift.
		* Esta es la logica de negocio real; el stub/skeleton
		* son transparentes para el programador.
		*/
		public class CalculadoraImpl implements Calculadora.Iface {
			
			@Override
			public double sumar(double a, double b) throws TException {
				System.out.printf("RPC sumar(%.2f, %.2f)%n", a, b);
				return a + b;
			}
			
			@Override
			public double restar(double a, double b) throws TException {
				System.out.printf("RPC restar(%.2f, %.2f)%n", a, b);
				return a - b;
			}
			
			@Override
			public double multiplicar(double a, double b) throws TException {
				System.out.printf("RPC multiplicar(%.2f, %.2f)%n", a, b);
				return a * b;
			}
			
			@Override
			public double dividir(double a, double b)
			throws OperacionInvalidaException, TException {
				System.out.printf("RPC dividir(%.2f, %.2f)%n", a, b);
				if (b == 0.0) {
					throw new OperacionInvalidaException(
					"Division por cero no permitida");
				}
				return a / b;
			}
			
			@Override
			public String nombreServidor() throws TException {
				return "ServidorThrift@" + java.net.InetAddress.getLocalHost().getHostName();
			}
		}
	\end{lstlisting}
	
	\paso{Crear \archivo{ServidorRPC.java} que arranca el servidor Thrift:}
	
	\begin{lstlisting}[style=JavaStyle, caption={ServidorRPC.java --- Servidor multihilo Thrift}]
		package ec.edu.uteq.distribuidas.rpc.server;
		
		import ec.edu.uteq.distribuidas.rpc.gen.Calculadora;
		import org.apache.thrift.server.*;
		import org.apache.thrift.transport.*;
		
		/**
		* Servidor RPC usando Apache Thrift.
		* TThreadPoolServer usa un pool de hilos para atender
		* multiples clientes en paralelo.
		*/
		public class ServidorRPC {
			
			private static final int PUERTO = 9090;
			
			public static void main(String[] args) throws Exception {
				Calculadora.Processor<CalculadoraImpl> processor =
				new Calculadora.Processor<>(new CalculadoraImpl());
				
				TServerTransport transport = new TServerSocket(PUERTO);
				
				TThreadPoolServer.Args serverArgs =
				new TThreadPoolServer.Args(transport)
				.processor(processor)
				.minWorkerThreads(5)
				.maxWorkerThreads(20);
				
				TServer server = new TThreadPoolServer(serverArgs);
				System.out.println("Servidor Thrift (RPC) escuchando en :" + PUERTO);
				server.serve();   // BLOQUEA hasta que el servidor se detenga
			}
		}
	\end{lstlisting}
	
	\paso{Crear \archivo{ClienteRPC.java}:}
	
	\begin{lstlisting}[style=JavaStyle, caption={ClienteRPC.java --- Cliente RPC con Thrift}]
		package ec.edu.uteq.distribuidas.rpc.client;
		
		import ec.edu.uteq.distribuidas.rpc.gen.*;
		import org.apache.thrift.protocol.*;
		import org.apache.thrift.transport.*;
		
		/**
		* Cliente que invoca metodos remotos en el servidor Thrift.
		* Desde el codigo del cliente, llamar dividir() parece una
		* llamada local, pero en realidad viaja por la red (RPC).
		*/
		public class ClienteRPC {
			
			private static final String HOST  = "localhost";
			private static final int    PUERTO = 9090;
			
			public static void main(String[] args) throws Exception {
				try (TTransport transport = new TSocket(HOST, PUERTO)) {
					transport.open();
					
					TProtocol protocol = new TBinaryProtocol(transport);
					Calculadora.Client cliente = new Calculadora.Client(protocol);
					
					System.out.println("Conectado al servidor: " + cliente.nombreServidor());
					System.out.printf("5 + 3 = %.2f%n", cliente.sumar(5, 3));
					System.out.printf("10 - 4 = %.2f%n", cliente.restar(10, 4));
					System.out.printf("6 x 7 = %.2f%n", cliente.multiplicar(6, 7));
					System.out.printf("15 / 4 = %.4f%n", cliente.dividir(15, 4));
					
					// Prueba de excepcion remota
					try {
						cliente.dividir(10, 0);
					} catch (OperacionInvalidaException ex) {
						System.out.println("Excepcion capturada: " + ex.getMessage());
					}
				}
			}
		}
	\end{lstlisting}
	
	\section{Ejecución}
	
	\paso{Ejecutar \textbf{ServidorRPC} primero. Salida esperada:}
	
	\begin{resultadobox}
		\begin{lstlisting}[style=BashStyle]
			Servidor Thrift (RPC) escuchando en :9090
		\end{lstlisting}
	\end{resultadobox}
	
	\paso{Ejecutar \textbf{ClienteRPC}. Salida esperada:}
	
	\begin{resultadobox}
		\begin{lstlisting}[style=BashStyle]
			Conectado al servidor: ServidorThrift@mi-pc
			5 + 3 = 8.00
			10 - 4 = 6.00
			6 x 7 = 42.00
			15 / 4 = 3.7500
			Excepcion capturada: Division por cero no permitida
		\end{lstlisting}
	\end{resultadobox}
	
	% ════════════════════════════════════════════════════════════
	%  GUÍA 4 — gRPC
	% ════════════════════════════════════════════════════════════
	\chapter{Guía 4: gRPC con Protocol Buffers}
	\resetpasos
	
	\begin{objetivobox}{gRPC con Protocol Buffers}
		Implementar un sistema de comunicación de alto rendimiento usando gRPC con
		Protocol Buffers (Protobuf) sobre HTTP/2, incluyendo llamadas unarias y
		streaming bidireccional~\parencite{indrasiri2020}.
	\end{objetivobox}
	
	\section{Fundamento teórico}
	
	\textbf{gRPC} (Google Remote Procedure Call) es un framework RPC moderno
	que usa \textbf{HTTP/2} como transporte y \textbf{Protocol Buffers} (Protobuf)
	como formato de serialización binario~\parencite{indrasiri2020}.
	
	Ventajas sobre alternativas:
	\begin{itemize}
		\item \textbf{HTTP/2}: multiplexing de múltiples RPCs sobre una sola conexión
		TCP (sin head-of-line blocking), compresión de headers (HPACK).
		\item \textbf{Protobuf}: 3-10$\times$ más compacto que JSON; 5-10$\times$ más
		rápido de serializar/deserializar.
		\item \textbf{Tipado fuerte}: el compilador \code{protoc} genera código
		verificado en tiempo de compilación para más de 20 lenguajes.
		\item \textbf{Streaming bidireccional}: cliente y servidor pueden enviarse
		flujos de mensajes simultáneamente.
	\end{itemize}
	
	gRPC define cuatro tipos de llamadas:
	\begin{enumerate}
		\item \textbf{Unaria}: cliente envía un mensaje, recibe uno.
		\item \textbf{Server streaming}: cliente envía uno, recibe un stream.
		\item \textbf{Client streaming}: cliente envía stream, recibe uno.
		\item \textbf{Bidireccional}: ambos envían y reciben streams simultáneamente.
	\end{enumerate}
	
	\section{Configuración del proyecto}
	
	\paso{Crear proyecto Maven en IntelliJ con ArtifactId \code{guia04-grpc}.
		Configurar \archivo{pom.xml} con las dependencias de gRPC y el plugin protoc:}
	
	\begin{lstlisting}[style=XMLStyle, caption={pom.xml --- Dependencias gRPC con generación automática}]
		<properties>
		<maven.compiler.source>21</maven.compiler.source>
		<maven.compiler.target>21</maven.compiler.target>
		<project.build.sourceEncoding>UTF-8</project.build.sourceEncoding>
		<grpc.version>1.62.2</grpc.version>
		<protobuf.version>3.25.3</protobuf.version>
		<protoc.version>3.25.3</protoc.version>
		</properties>
		
		<dependencies>
		<dependency>
		<groupId>io.grpc</groupId>
		<artifactId>grpc-netty-shaded</artifactId>
		<version>${grpc.version}</version>
		</dependency>
		<dependency>
		<groupId>io.grpc</groupId>
		<artifactId>grpc-protobuf</artifactId>
		<version>${grpc.version}</version>
		</dependency>
		<dependency>
		<groupId>io.grpc</groupId>
		<artifactId>grpc-stub</artifactId>
		<version>${grpc.version}</version>
		</dependency>
		<dependency>
		<groupId>com.google.protobuf</groupId>
		<artifactId>protobuf-java</artifactId>
		<version>${protobuf.version}</version>
		</dependency>
		<!-- Anotacion @Generated requerida por el codigo generado -->
		<dependency>
		<groupId>javax.annotation</groupId>
		<artifactId>javax.annotation-api</artifactId>
		<version>1.3.2</version>
		</dependency>
		</dependencies>
		
		<build>
		<extensions>
		<extension>
		<groupId>kr.motd.maven</groupId>
		<artifactId>os-maven-plugin</artifactId>
		<version>1.7.1</version>
		</extension>
		</extensions>
		<plugins>
		<plugin>
		<groupId>org.xolstice.maven.plugins</groupId>
		<artifactId>protobuf-maven-plugin</artifactId>
		<version>0.6.1</version>
		<configuration>
		<protocArtifact>
		com.google.protobuf:protoc:${protoc.version}:exe:${os.detected.classifier}
		</protocArtifact>
		<pluginId>grpc-java</pluginId>
		<pluginArtifact>
		io.grpc:protoc-gen-grpc-java:${grpc.version}:exe:${os.detected.classifier}
		</pluginArtifact>
		</configuration>
		<executions>
		<execution>
		<goals>
		<goal>compile</goal>
		<goal>compile-custom</goal>
		</goals>
		</execution>
		</executions>
		</plugin>
		</plugins>
		</build>
	\end{lstlisting}
	
	\section{Definición del contrato Protobuf}
	
	\paso{Crear el directorio \code{src/main/proto} y el archivo
		\archivo{mensajeria.proto}:}
	
	\begin{lstlisting}[style=ProtoStyle, caption={mensajeria.proto --- Definición del servicio gRPC}]
		syntax = "proto3";
		
		package ec.edu.uteq.distribuidas.grpc;
		option java_package = "ec.edu.uteq.distribuidas.grpc.gen";
		option java_outer_classname = "MensajeriaProto";
		option java_multiple_files = true;
		
		// -- Mensajes de solicitud y respuesta -------------
		message SolicitudMensaje {
			string remitente   = 1;
			string contenido   = 2;
			int64  timestamp   = 3;  // Reloj de Lamport
		}
		
		message RespuestaMensaje {
			string servidor    = 1;
			string contenido   = 2;
			int64  timestamp   = 3;
			bool   exito       = 4;
		}
		
		message SolicitudStream {
			string usuario = 1;
			int32  cantidad = 2;  // Cuantos mensajes transmitir
		}
		
		// -- Definicion del servicio gRPC -------------------
		service ServicioMensajeria {
			
			// Llamada unaria: cliente envia 1 mensaje, recibe 1 respuesta
			rpc EnviarMensaje(SolicitudMensaje) returns (RespuestaMensaje);
			
			// Server streaming: el servidor envia multiples respuestas
			rpc RecibirNoticias(SolicitudStream) returns (stream RespuestaMensaje);
			
			// Bidireccional: cliente y servidor intercambian streams
			rpc ChatBidireccional(stream SolicitudMensaje)
			returns (stream RespuestaMensaje);
		}
	\end{lstlisting}
	
	\paso{Generar el código Java ejecutando desde la terminal de IntelliJ:}
	
	\begin{lstlisting}[style=BashStyle]
		mvn generate-sources
	\end{lstlisting}
	
	El código se genera en \code{target/generated-sources/protobuf/}.
	IntelliJ lo detecta automáticamente si se ejecutó \menu{Maven $\to$ Reload Project}.
	
	\section{Implementación del servidor gRPC}
	
	\paso{Crear \archivo{MensajeriaServiceImpl.java}:}
	
	\begin{lstlisting}[style=JavaStyle, caption={MensajeriaServiceImpl.java --- Implementación de los tres tipos de llamada}]
		package ec.edu.uteq.distribuidas.grpc.server;
		
		import ec.edu.uteq.distribuidas.grpc.gen.*;
		import io.grpc.stub.StreamObserver;
		
		import java.util.concurrent.atomic.AtomicLong;
		
		/**
		* Implementacion de los tres tipos de llamada gRPC:
		*   1. Unaria
		*   2. Server streaming
		*   3. Bidireccional
		*/
		public class MensajeriaServiceImpl
		extends ServicioMensajeriaGrpc.ServicioMensajeriaImplBase {
			
			private final AtomicLong relojLamport = new AtomicLong(0);
			
			// -- 1. Llamada unaria -----------------------------
			@Override
			public void enviarMensaje(SolicitudMensaje request,
			StreamObserver<RespuestaMensaje> responseObserver) {
				// Regla R3 del reloj de Lamport
				long ts = Math.max(relojLamport.get(), request.getTimestamp()) + 1;
				relojLamport.set(ts);
				
				System.out.printf("[gRPC Unario] %s: %s (ts=%d)%n",
				request.getRemitente(), request.getContenido(), ts);
				
				RespuestaMensaje respuesta = RespuestaMensaje.newBuilder()
				.setServidor("ServidorGRPC")
				.setContenido("ACK: " + request.getContenido())
				.setTimestamp(ts)
				.setExito(true)
				.build();
				
				responseObserver.onNext(respuesta);
				responseObserver.onCompleted();
			}
			
			// -- 2. Server streaming ---------------------------
			@Override
			public void recibirNoticias(SolicitudStream request,
			StreamObserver<RespuestaMensaje> responseObserver) {
				String usuario = request.getUsuario();
				int cantidad   = request.getCantidad();
				
				System.out.printf("[gRPC Streaming] Enviando %d noticias a %s%n",
				cantidad, usuario);
				
				for (int i = 1; i <= cantidad; i++) {
					long ts = relojLamport.incrementAndGet();
					RespuestaMensaje noticia = RespuestaMensaje.newBuilder()
					.setServidor("ServidorGRPC")
					.setContenido("Noticia #" + i + " para " + usuario)
					.setTimestamp(ts)
					.setExito(true)
					.build();
					responseObserver.onNext(noticia);
					
					try { Thread.sleep(500); } catch (InterruptedException ignored) {}
				}
				responseObserver.onCompleted();
			}
			
			// -- 3. Chat bidireccional -------------------------
			@Override
			public StreamObserver<SolicitudMensaje> chatBidireccional(
			StreamObserver<RespuestaMensaje> responseObserver) {
				
				return new StreamObserver<>() {
					@Override
					public void onNext(SolicitudMensaje msg) {
						long ts = Math.max(relojLamport.get(), msg.getTimestamp()) + 1;
						relojLamport.set(ts);
						System.out.printf("[gRPC Bidir] %s: %s%n",
						msg.getRemitente(), msg.getContenido());
						// Eco con el timestamp actualizado
						responseObserver.onNext(RespuestaMensaje.newBuilder()
						.setServidor("ServidorGRPC")
						.setContenido("[" + msg.getRemitente() + "] "
						+ msg.getContenido())
						.setTimestamp(ts)
						.setExito(true)
						.build());
					}
					
					@Override
					public void onError(Throwable t) {
						System.err.println("Error en stream: " + t.getMessage());
					}
					
					@Override
					public void onCompleted() {
						responseObserver.onCompleted();
					}
				};
			}
		}
	\end{lstlisting}
	
	\paso{Crear \archivo{ServidorGRPC.java}:}
	
	\begin{lstlisting}[style=JavaStyle, caption={ServidorGRPC.java --- Servidor gRPC sobre HTTP/2}]
		package ec.edu.uteq.distribuidas.grpc.server;
		
		import io.grpc.Server;
		import io.grpc.ServerBuilder;
		
		public class ServidorGRPC {
			
			private static final int PUERTO = 50051;
			
			public static void main(String[] args) throws Exception {
				Server server = ServerBuilder.forPort(PUERTO)
				.addService(new MensajeriaServiceImpl())
				.build()
				.start();
				
				System.out.println("Servidor gRPC iniciado en puerto " + PUERTO);
				Runtime.getRuntime().addShutdownHook(new Thread(server::shutdown));
				server.awaitTermination();
			}
		}
	\end{lstlisting}
	
	\paso{Crear \archivo{ClienteGRPC.java} para probar los tres tipos de llamada:}
	
	\begin{lstlisting}[style=JavaStyle, caption={ClienteGRPC.java --- Cliente que demuestra los tres tipos de llamada}]
		package ec.edu.uteq.distribuidas.grpc.client;
		
		import ec.edu.uteq.distribuidas.grpc.gen.*;
		import io.grpc.*;
		import io.grpc.stub.StreamObserver;
		
		import java.util.Iterator;
		import java.util.concurrent.*;
		import java.util.concurrent.atomic.AtomicLong;
		
		public class ClienteGRPC {
			
			private static final AtomicLong reloj = new AtomicLong(0);
			
			public static void main(String[] args) throws Exception {
				ManagedChannel canal = ManagedChannelBuilder
				.forAddress("localhost", 50051)
				.usePlaintext()           // Sin TLS (solo para desarrollo)
				.build();
				
				ServicioMensajeriaGrpc.ServicioMensajeriaBlockingStub
				stubBloqueante = ServicioMensajeriaGrpc.newBlockingStub(canal);
				ServicioMensajeriaGrpc.ServicioMensajeriaStub
				stubAsync = ServicioMensajeriaGrpc.newStub(canal);
				
				// -- 1. Llamada unaria -------------------------
				System.out.println("=== 1. Llamada Unaria ===");
				SolicitudMensaje solicitud = SolicitudMensaje.newBuilder()
				.setRemitente("ClienteA")
				.setContenido("Hola desde gRPC")
				.setTimestamp(reloj.incrementAndGet())
				.build();
				RespuestaMensaje resp = stubBloqueante.enviarMensaje(solicitud);
				System.out.println("Respuesta: " + resp.getContenido()
				+ " (ts=" + resp.getTimestamp() + ")");
				
				// -- 2. Server streaming -----------------------
				System.out.println("\n=== 2. Server Streaming (3 noticias) ===");
				Iterator<RespuestaMensaje> noticias = stubBloqueante.recibirNoticias(
				SolicitudStream.newBuilder()
				.setUsuario("ClienteA")
				.setCantidad(3)
				.build());
				while (noticias.hasNext()) {
					System.out.println("  << " + noticias.next().getContenido());
				}
				
				// -- 3. Bidireccional -------------------------
				System.out.println("\n=== 3. Chat Bidireccional ===");
				CountDownLatch latch = new CountDownLatch(1);
				StreamObserver<SolicitudMensaje> emisor =
				stubAsync.chatBidireccional(new StreamObserver<>() {
					@Override
					public void onNext(RespuestaMensaje r) {
						System.out.println("  << " + r.getContenido());
					}
					@Override
					public void onError(Throwable t) {
						latch.countDown();
					}
					@Override
					public void onCompleted() {
						System.out.println("Stream completado.");
						latch.countDown();
					}
				});
				
				// Enviar 3 mensajes en el stream del cliente
				for (int i = 1; i <= 3; i++) {
					emisor.onNext(SolicitudMensaje.newBuilder()
					.setRemitente("ClienteA")
					.setContenido("Mensaje bidireccional #" + i)
					.setTimestamp(reloj.incrementAndGet())
					.build());
					Thread.sleep(300);
				}
				emisor.onCompleted();
				latch.await(5, TimeUnit.SECONDS);
				
				canal.shutdown();
			}
		}
	\end{lstlisting}
	
	\section{Ejecución}
	
	\paso{Ejecutar \menu{Maven $\to$ generate-sources} para compilar el \code{.proto}.
		Luego ejecutar \textbf{ServidorGRPC}. Verificar:}
	
	\begin{resultadobox}
		\begin{lstlisting}[style=BashStyle]
			Servidor gRPC iniciado en puerto 50051
		\end{lstlisting}
	\end{resultadobox}
	
	\paso{Ejecutar \textbf{ClienteGRPC}. Salida esperada:}
	
	\begin{resultadobox}
		\begin{lstlisting}[style=BashStyle]
			=== 1. Llamada Unaria ===
			Respuesta: ACK: Hola desde gRPC (ts=2)
			
			=== 2. Server Streaming (3 noticias) ===
			<< Noticia #1 para ClienteA
			<< Noticia #2 para ClienteA
			<< Noticia #3 para ClienteA
			
			=== 3. Chat Bidireccional ===
			<< [ClienteA] Mensaje bidireccional #1
			<< [ClienteA] Mensaje bidireccional #2
			<< [ClienteA] Mensaje bidireccional #3
			Stream completado.
		\end{lstlisting}
	\end{resultadobox}
	
	% ════════════════════════════════════════════════════════════
	%  GUÍA 5 — RMI
	% ════════════════════════════════════════════════════════════
	\chapter{Guía 5: Java RMI (Remote Method Invocation)}
	\resetpasos
	
	\begin{objetivobox}{Java RMI}
		Implementar un sistema RMI nativo de Java que permita invocar métodos de objetos
		remotos como si fueran locales, usando el registro RMI como servicio de nombres
		y la serialización de objetos Java~\parencite{waldo1994}.
	\end{objetivobox}
	
	\section{Fundamento teórico}
	
	\textbf{RMI} (Remote Method Invocation) es la implementación de RPC de Java,
	introducida en Java 1.1. Permite a un programa Java invocar métodos de un objeto
	Java ubicado en otra JVM (Java Virtual Machine), posiblemente en otro host~\parencite{waldo1994}.
	
	Arquitectura de tres capas:
	\begin{itemize}
		\item \textbf{Capa de stub/skeleton}: generados en tiempo de compilación
		(Java 5+: generados automáticamente, no requiere \code{rmic}).
		\item \textbf{Capa de referencia remota}: maneja la semántica de llamada
		(at-most-once), gestión de la conexión TCP subyacente.
		\item \textbf{Capa de transporte}: usa JRMP (\textit{Java Remote Method Protocol})
		sobre TCP. Serialización nativa de Java.
	\end{itemize}
	
	El \textbf{RMI Registry} (puerto 1099 por defecto) actúa como servicio de
	nombres: el servidor registra sus objetos con \code{Naming.bind()},
	el cliente los localiza con \code{Naming.lookup()}~\parencite{vanSteen2023}.
	
	\begin{figure}[H]
		\centering
		\begin{tikzpicture}[node distance=0.5cm and 2.5cm, font=\small]
			% JVM Cliente
			\node[clibox] (app)    {App Cliente};
			\node[netbox, below=0.8cm of app] (stub)  {Stub (proxy)};
			\draw[-Stealth, thick] (app) -- node[right, font=\tiny]{invoca} (stub);
			
			% RMI Registry
			\node[netbox, right=2.5cm of app] (reg)   {RMI Registry\\(puerto 1099)};
			
			% JVM Servidor
			\node[srvbox, right=2.5cm of reg] (impl)  {Impl. remota};
			\node[netbox, below=0.8cm of impl] (skel) {Skeleton};
			\draw[-Stealth, thick] (skel) -- node[right, font=\tiny]{delega} (impl);
			
			% Flujos
			\draw[msgArrow] (app.east) -- node[above, font=\tiny]{lookup(''Calculadora'')}
			(reg.west);
			\draw[retArrow] (reg.west) -- node[below, font=\tiny]{stub remoto}
			([yshift=-6pt]app.east);
			\draw[msgArrow] (stub.east) -- node[above, font=\tiny]{JRMP/TCP}
			(skel.west);
			\draw[retArrow] (skel.west) -- node[below, font=\tiny]{resultado}
			([yshift=-6pt]stub.east);
			\draw[msgArrow] (impl.west) to[bend left=30]
			node[above, font=\tiny]{bind(''Calculadora'')} (reg.east);
			
			% Etiquetas JVM
			\node[above=0.1cm of app, font=\bfseries\small,
			text=practiceBlue] {JVM Cliente};
			\node[above=0.1cm of impl, font=\bfseries\small,
			text=uteqDark] {JVM Servidor};
			\node[above=0.1cm of reg, font=\bfseries\small,
			text=gray!70] {Registry};
		\end{tikzpicture}
		\caption{Arquitectura completa del sistema RMI con tres JVMs}
		\label{fig:rmi_arquitectura}
	\end{figure}
	
	\section{Estructura del proyecto}
	
	\paso{Crear proyecto Maven con ArtifactId \code{guia05-rmi}.
		No se requieren dependencias externas: RMI es parte del JDK estándar.}
	
	\begin{lstlisting}[style=XMLStyle, caption={pom.xml para proyecto RMI}]
		<properties>
		<maven.compiler.source>21</maven.compiler.source>
		<maven.compiler.target>21</maven.compiler.target>
		<project.build.sourceEncoding>UTF-8</project.build.sourceEncoding>
		</properties>
		<!-- Sin dependencias adicionales: RMI esta en java.rmi del JDK -->
	\end{lstlisting}
	
	\section{Interfaz remota}
	
	\paso{Crear el paquete \code{ec.edu.uteq.distribuidas.rmi.api} y la
		interfaz remota \archivo{CalculadoraRemota.java}:}
	
	\begin{lstlisting}[style=JavaStyle, caption={CalculadoraRemota.java --- Interfaz del objeto remoto}]
		package ec.edu.uteq.distribuidas.rmi.api;
		
		import java.rmi.Remote;
		import java.rmi.RemoteException;
		
		/**
		* Interfaz remota RMI.
		* REGLAS OBLIGATORIAS para interfaces RMI:
		*   1. Debe extender java.rmi.Remote.
		*   2. Cada metodo DEBE declarar throws RemoteException
		*      (para manejar fallos de red transparentemente).
		*   3. Los parametros y el valor de retorno deben ser:
		*      - Tipos primitivos, o
		*      - Implementar java.io.Serializable, o
		*      - Implementar Remote (objetos remotos anidados).
		*/
		public interface CalculadoraRemota extends Remote {
			
			double sumar(double a, double b) throws RemoteException;
			double restar(double a, double b) throws RemoteException;
			double multiplicar(double a, double b) throws RemoteException;
			double dividir(double a, double b) throws RemoteException;
			String obtenerInfo() throws RemoteException;
			
			// Metodos adicionales (Java 21)
			double potencia(double base, double exponente) throws RemoteException;
			double raizCuadrada(double n) throws RemoteException;
		}
	\end{lstlisting}
	
	\section{Implementación del objeto remoto}
	
	\paso{Crear el paquete \code{ec.edu.uteq.distribuidas.rmi.server} y la implementación \archivo{CalculadoraImpl.java}:}
	
	\begin{lstlisting}[style=JavaStyle, caption={CalculadoraImpl.java --- Implementación del objeto remoto}]
		package ec.edu.uteq.distribuidas.rmi.server;
		
		import ec.edu.uteq.distribuidas.rmi.api.CalculadoraRemota;
		import java.rmi.RemoteException;
		import java.rmi.server.UnicastRemoteObject;
		
		/**
		* Implementacion concreta del objeto remoto.
		*
		* OBLIGATORIO: extender UnicastRemoteObject.
		* Esto hace que el objeto:
		*   - Se exporte automaticamente al instanciarse.
		*   - Use TCP para recibir llamadas remotas.
		*   - Sea elegible para garbage collection remoto.
		*
		* Alternativa: NO extender UnicastRemoteObject y usar
		*   UnicastRemoteObject.exportObject(this, 0) manualmente.
		*/
		public class CalculadoraImpl
		extends UnicastRemoteObject
		implements CalculadoraRemota {
			
			private static final long serialVersionUID = 1L;
			
			private final String nombreServidor;
			
			// El constructor DEBE declarar throws RemoteException
			public CalculadoraImpl(String nombre) throws RemoteException {
				super();   // exporta el objeto (asigna puerto TCP aleatorio)
				this.nombreServidor = nombre;
			}
			
			@Override
			public double sumar(double a, double b) throws RemoteException {
				logLlamada("sumar", a, b);
				return a + b;
			}
			
			@Override
			public double restar(double a, double b) throws RemoteException {
				logLlamada("restar", a, b);
				return a - b;
			}
			
			@Override
			public double multiplicar(double a, double b) throws RemoteException {
				logLlamada("multiplicar", a, b);
				return a * b;
			}
			
			@Override
			public double dividir(double a, double b) throws RemoteException {
				logLlamada("dividir", a, b);
				if (b == 0.0) {
					// Se puede lanzar RemoteException o una excepcion serializable propia
					throw new RemoteException("Division por cero no permitida");
				}
				return a / b;
			}
			
			@Override
			public double potencia(double base, double exp) throws RemoteException {
				return Math.pow(base, exp);
			}
			
			@Override
			public double raizCuadrada(double n) throws RemoteException {
				if (n < 0) throw new RemoteException("Raiz de negativo no definida en R");
				return Math.sqrt(n);
			}
			
			@Override
			public String obtenerInfo() throws RemoteException {
				return "Servidor RMI: " + nombreServidor
				+ " | JVM: " + System.getProperty("java.version")
				+ " | Hilos activos: " + Thread.activeCount();
			}
			
			private void logLlamada(String metodo, double a, double b) {
				System.out.printf("[RMI] %s(%.2f, %.2f) invocado%n", metodo, a, b);
			}
		}
	\end{lstlisting}
	
	\paso{Crear \archivo{ServidorRMI.java} que arranca el registry y exporta el objeto:}
	
	\begin{lstlisting}[style=JavaStyle, caption={ServidorRMI.java --- Servidor que registra el objeto remoto}]
		package ec.edu.uteq.distribuidas.rmi.server;
		
		import java.rmi.Naming;
		import java.rmi.registry.LocateRegistry;
		
		/**
		* Servidor RMI:
		*   1. Crea el RMI Registry en el puerto 1099.
		*   2. Instancia el objeto remoto (lo exporta automaticamente).
		*   3. Lo registra en el Registry con un nombre.
		*/
		public class ServidorRMI {
			
			private static final int PUERTO_REGISTRY = 1099;
			private static final String NOMBRE       = "CalculadoraDistribuida";
			
			public static void main(String[] args) {
				try {
					// Paso 1: Crear el registry en este proceso
					LocateRegistry.createRegistry(PUERTO_REGISTRY);
					System.out.println("RMI Registry creado en puerto " + PUERTO_REGISTRY);
					
					// Paso 2: Instanciar e implementacion (se exporta automaticamente)
					CalculadoraImpl calculadora = new CalculadoraImpl("UTEQ-Servidor-01");
					
					// Paso 3: Registrar con el nombre buscable por clientes
					// URL: rmi://host:puerto/nombre
					String url = "rmi://localhost:" + PUERTO_REGISTRY + "/" + NOMBRE;
					Naming.bind(url, calculadora);
					
					System.out.println("Objeto registrado en: " + url);
					System.out.println("Servidor RMI listo. Esperando clientes...");
					System.out.println(calculadora.obtenerInfo());
					
					// El servidor se mantiene vivo mientras el registry exista
					// En produccion usar join() en lugar de sleep
					Thread.currentThread().join();
					
				} catch (Exception e) {
					System.err.println("Error en ServidorRMI: " + e.getMessage());
					e.printStackTrace();
				}
			}
		}
	\end{lstlisting}
	
	\section{Cliente RMI}
	
	\paso{Crear el paquete \code{ec.edu.uteq.distribuidas.rmi.client} y
		la clase \archivo{ClienteRMI.java}:}
	
	\begin{lstlisting}[style=JavaStyle, caption={ClienteRMI.java --- Cliente que invoca métodos remotos}]
		package ec.edu.uteq.distribuidas.rmi.client;
		
		import ec.edu.uteq.distribuidas.rmi.api.CalculadoraRemota;
		import java.rmi.Naming;
		import java.rmi.RemoteException;
		
		/**
		* Cliente RMI.
		* IMPORTANTE: El cliente SOLO necesita la interfaz CalculadoraRemota.java.
		* NO necesita la implementacion CalculadoraImpl.java.
		* Esta es la "transparencia de acceso" de RMI.
		*/
		public class ClienteRMI {
			
			private static final String HOST   = "localhost";
			private static final int    PUERTO = 1099;
			private static final String NOMBRE = "CalculadoraDistribuida";
			
			public static void main(String[] args) {
				String url = "rmi://" + HOST + ":" + PUERTO + "/" + NOMBRE;
				
				try {
					// lookup() devuelve un stub (proxy) que parece un objeto local
					CalculadoraRemota calc =
					(CalculadoraRemota) Naming.lookup(url);
					
					System.out.println("Conectado a: " + calc.obtenerInfo());
					System.out.println("-----------------------------------");
					
					// Todas estas llamadas viajan por la red transparentemente
					System.out.printf("10 + 5   = %.2f%n", calc.sumar(10, 5));
					System.out.printf("10 - 5   = %.2f%n", calc.restar(10, 5));
					System.out.printf("10 x 5   = %.2f%n", calc.multiplicar(10, 5));
					System.out.printf("10 / 4   = %.4f%n", calc.dividir(10, 4));
					System.out.printf("2 ^ 10   = %.0f%n", calc.potencia(2, 10));
					System.out.printf("sqrt(2)  = %.6f%n", calc.raizCuadrada(2));
					
					// Prueba de excepcion remota
					try {
						calc.dividir(5, 0);
					} catch (RemoteException e) {
						System.out.println("Excepcion remota: " + e.getCause().getMessage());
					}
					
				} catch (Exception e) {
					System.err.println("Error en ClienteRMI: " + e.getMessage());
				}
			}
		}
	\end{lstlisting}
	
	\section{Ejecución}
	
	\paso{Ejecutar \textbf{ServidorRMI} primero. Salida esperada:}
	
	\begin{resultadobox}
		\begin{lstlisting}[style=BashStyle]
			RMI Registry creado en puerto 1099
			Objeto registrado en: rmi://localhost:1099/CalculadoraDistribuida
			Servidor RMI listo. Esperando clientes...
			Servidor RMI: UTEQ-Servidor-01 | JVM: 21.0.3 | Hilos activos: 4
		\end{lstlisting}
	\end{resultadobox}
	
	\paso{Ejecutar \textbf{ClienteRMI}. Salida esperada:}
	
	\begin{resultadobox}
		\begin{lstlisting}[style=BashStyle]
			Conectado a: Servidor RMI: UTEQ-Servidor-01 | JVM: 21.0.3 | Hilos activos: 5
			-----------------------------------
			10 + 5   = 15.00
			10 - 5   = 5.00
			10 x 5   = 50.00
			10 / 4   = 2.5000
			2 ^ 10   = 1024
			sqrt(2)  = 1.414214
			Excepcion remota: Division por cero no permitida
		\end{lstlisting}
	\end{resultadobox}
	
	\begin{advertenciabox}
		Si aparece \code{java.rmi.ConnectException: Connection refused to host}: verificar que
		el \textbf{ServidorRMI} esté ejecutándose antes que el cliente. En redes locales con
		múltiples interfaces de red, puede ser necesario agregar
		\code{-Djava.rmi.server.hostname=IP\_del\_servidor} como argumento de JVM al servidor.
	\end{advertenciabox}
	
	\paso{Para simular dos JVMs en equipos distintos: cambiar \code{HOST} en
		\archivo{ClienteRMI.java} a la IP del servidor. El cliente solo necesita la
		interfaz \code{CalculadoraRemota.java} en su classpath.}
	
	% ════════════════════════════════════════════════════════════
	%  ANÁLISIS COMPARATIVO
	% ════════════════════════════════════════════════════════════
	\chapter{Análisis Comparativo}
	
	\section{Cuadro sinóptico de mecanismos}
	
	\begin{table}[H]
		\centering
		\caption{Comparativa técnica de los cinco mecanismos de comunicación implementados}
		\label{tab:comparativa_final}
		\setlength{\tabcolsep}{4pt}
		\begin{tabularx}{\linewidth}{@{}lXXXXX@{}}
			\toprule
			\textbf{Criterio} &
			\textbf{TCP/UDP} &
			\textbf{WebSocket} &
			\textbf{Thrift RPC} &
			\textbf{gRPC} &
			\textbf{RMI} \\
			\midrule
			Protocolo base    & TCP / UDP  & HTTP $\to$ WS & TCP (bin)   & HTTP/2     & JRMP/TCP \\
			Serialización     & Manual     & JSON/texto     & Thrift bin  & Protobuf   & Java ser.\\
			Tipado IDL        & Ninguno    & Ninguno        & \code{.thrift} & \code{.proto} & Java iface\\
			Streaming         & Sí (TCP)   & Bidireccional  & No nativo   & Sí (4 tipos)& No     \\
			Multilenguaje     & Sí         & Sí             & 20+ lang.   & 20+ lang.  & Solo Java\\
			Overhead (aprox.) & Muy bajo   & Bajo           & Bajo        & Muy bajo   & Alto     \\
			Complejidad conf. & Baja       & Media          & Media       & Media      & Baja     \\
			Caso de uso ideal & Chat, bin  & Tiempo real    & Microserv.  & Microserv. & Java-Java\\
			\bottomrule
		\end{tabularx}
	\end{table}
	
	\section{Criterios de selección}
	
	Para elegir el mecanismo de comunicación en un proyecto distribuido, considerar:
	
	\begin{enumerate}
		\item \textbf{¿Es solo Java?} RMI es la opción más simple. De lo contrario,
		descartar RMI.
		\item \textbf{¿Requiere tiempo real en el browser?} WebSocket es la única
		opción con soporte nativo en navegadores sin polling.
		\item \textbf{¿Rendimiento crítico entre microservicios?} gRPC con Protobuf
		ofrece la mejor relación latencia/throughput.
		\item \textbf{¿Control total sobre la red?} Sockets TCP/UDP permiten optimizar
		cada byte, al costo de mayor complejidad de desarrollo.
		\item \textbf{¿Necesita interoperar con sistemas legados?} Apache Thrift tiene
		soporte para más lenguajes incluyendo C++ y Erlang.
	\end{enumerate}
	
	% ════════════════════════════════════════════════════════════
	%  BIBLIOGRAFÍA
	% ════════════════════════════════════════════════════════════
	\printbibliography[title={Referencias Bibliográficas}]
	
\end{document}